\section{Métricas de sucesso e experimentos planejados}
\label{sec:metricas}

\begin{description}
  \item[\emph{Speedup}:] $S = T_1 / T_N$, razão entre o tempo com 1 worker
    e com $N$ workers. O ideal seria $S = N$, mas a parte que não se
    paraleliza (ler os dados, agregar, comunicar) limita o ganho (Lei de
    Amdahl).
  \item[Eficiência:] $E = S / N$, quanto de cada worker adicional é
    aproveitado.
  \item[Distribuição da carga:] pedaços processados e tempo ocioso de
    cada worker.
  \item[Corretude:] o resultado distribuído deve ser idêntico ao de uma
    execução sequencial de referência, inclusive quando um worker falha.
\end{description}

Experimentos planejados:

\begin{enumerate}
  \item \textbf{Escalabilidade:} tempo e \emph{speedup} com 1, 2 e 4
    workers, e o efeito do tamanho do pedaço, para identificar quando a
    comunicação passa a dominar.
  \item \textbf{Estática $\times$ dinâmica} em um cluster heterogêneo
    (máquinas diferentes ou atraso artificial em um worker): tempo total,
    pedaços por worker e ociosidade.
  \item \textbf{Tolerância a falhas:} derrubar um worker no meio da
    execução e verificar se o mestre reatribui o pedaço dele.
\end{enumerate}
